\begin{figure*}[t]
  \centering
  \blankfigure{22mm}{可视化结果待补}{输入 RGBA、二维 B/C/D 引导、对应三维结果与局部放大。使用同一对象、同一相机；同时保留成功与失败样本。}
  \caption{定性结果占位。计划展示组装结构、实体体积结构与细薄结构对象的二维引导及三维分割对比；本版不放入未经配对核验的结果。}
  \label{fig:qualitative}
\end{figure*}
\begin{figure*}[t]
  \centering
  \begin{minipage}[t]{0.48\textwidth}
    \centering
    \blankfigure{18mm}{分组性能柱状图待补}{按三类对象分别汇总主指标，附对象级配对置信区间。}
    \caption{结构类别对结果的影响（占位）。}
    \label{fig:group-bars}
  \end{minipage}\hfill
  \begin{minipage}[t]{0.48\textwidth}
    \centering
    \blankfigure{18mm}{消融差值柱状图待补}{以完整方法为参照，报告同一测试子集上的指标变化。}
    \caption{提示、反馈与融合的消融分析（占位）。}
    \label{fig:ablation-bars}
  \end{minipage}
\end{figure*}
\begin{figure*}[t]
  \centering
  \blankfigure{18mm}{提示预算—质量／耗时折线图待补}{横轴使用预先定义的点提示预算；分别绘制质量与端到端时间，不预设单调改善。}
  \caption{提示预算与质量、计算成本之间的关系（占位）。本图属于追加实验，不能从表\ref{tab:ablation}的若干异质配置直接连线得到。}
  \label{fig:budget-lines}
\end{figure*}
